\section{Steps cited without explanation}\label{sec:app-open}

The lists below name, section by section, the steps for which the text states the result and
cites the Lean text without saying why it holds, or says why only as far as the documentation of
the code does. Each item names the step and the place where it is used, cites the declarations or
files to read, and says what the text gives. Background results admitted as Facts are listed in
Appendix~\ref{sec:app-facts} and appear here only where the text describes in part how the Lean
text proves them; sections with no such steps have no list.

\begingroup\small

\subsection*{Section~\ref{sec:zeros}: From a one-sided bound to the zeros}
\begin{itemize}[leftmargin=2em,itemsep=1pt]
\item The prime number theorem for the remainder, Theorem~\ref{zeros:thm:pnt}: the zero-free
  region with the bound for $G_q(s)/s$ (\lean{New.eventually_norm_logDerivMean_div_le},
  \leanfile{Lib/Analytic/PrimeNumberTheorem/ZeroFreeRegion},
  \leanfile{Lib/Analytic/PrimeNumberTheorem/LFunctionGrowth}) and the Tauberian step
  (\lean{New.eventually_abs_le_inv_pow_of_transform}). The text describes the two steps in a
  paragraph and refers to the literature; the proof is not given.
\item Landau's lemma, Fact~\ref{zeros:fact:landau}. The text derives it from Landau's theorem on
  Laplace transforms of non-negative functions, with a reference; how the Lean text proves that
  theorem (\lean{New.mem_integrableExpSet_of_summable_taylor}, \leanfile{Lib/Analytic/Landau}) is
  not described.
\item The choice of the prime, Fact~\ref{zeros:fact:prime}, through the zero-density estimate for
  quadratic twists (\lean{New.zero_density_quadTwist}, \leanfile{Lib/Analytic/QuadraticTwist}). The
  text gives the role of the estimate, the sources, and the weak form proved in Lean in one
  sentence.
\end{itemize}

\subsection*{Section~\ref{sec:setup}: The data, prime sums, heights and twisted lattices}
\begin{itemize}[leftmargin=2em,itemsep=1pt]
\item The sizes of the data in Table~\ref{setup:tab:sizes}, and every inequality between sizes that
  the paper quotes by \lean{New.Data.le_of_hasLogBound} (for instance in the proofs of
  Lemma~\ref{data:lem:N} and of Theorems~\ref{main:thm:A} and~\ref{main:thm:star}). The text gives
  the method, a comparison of upper and lower bounds for logarithms, not the bounds themselves
  (\leanfile{New/Reduction/Data/Magnitudes}, \leanfile{New/Reduction/Data/Sizes}; for
  Section~\ref{sec:ub-divis} also \leanfile{New/UpperBound/Size/Estimates}).
\end{itemize}
The construction of the data behind Theorem~\ref{setup:thm:existence}, which
Section~\ref{sec:setup} only sketches, is carried out in Section~\ref{sec:data}; apart from the
comparisons of sizes in Lemma~\ref{data:lem:N}, that section leaves no step to list.

\subsection*{Section~\ref{sec:main}: Three theorems, the main inequality and the proof of the theorem}
\begin{itemize}[leftmargin=2em,itemsep=1pt]
\item The inequality \eqref{main:eq:cB} for $\log c_B$, used in the proof of
  Theorem~\ref{main:thm:A}(a) (\lean{New.Data.log_invGainFactor_le}). The text gives the statement
  and says only that the term $b$ bounds the loss from rounding the exponents.
\end{itemize}

\subsection*{Section~\ref{sec:ub-objects}: Theta kernels, source forms, target forms and chains of lifts}
\begin{itemize}[leftmargin=2em,itemsep=1pt]
\item The lift laws modulo $N$ (end of \S\ref{ub-objects:sec:composite}): $p$ glued double lifts of
  a ramified $\phi$ are a scalar times $2p$ direct lifts of $\mathrm{Up}^p\phi$.
  \lean{New.IsSplitLevel.compUpTo_gluedLift_eq_smul}. The text gives the statement only, not the laws
  of the mirror and direct data behind it.
\item The bound $\exp(x^{70})$ for the branching map (\S\ref{ub-objects:sec:composite}), used in the
  proof of Theorem~\ref{ub-divis:thm:size} through \lean{New.Data.weylNormSq_restrictToY_le}.
  \lean{New.Data.exists_sum_columnsY_le_and_pow_le_exp}. The text gives the statement only.
\item The gauge bound for standard data at $\ell$ not split in $K$, and the normalising constant of the
  definite theta kernels (after Proposition~\ref{ub-objects:prop:padic}).
  \lean{New.Parameters.mem_raises_predictedDivisibility}, \lean{New.Rows.exists_ne_zero_isModKernelScale}.
  The text gives the statements only; the count behind the first is Theorem~\ref{der:thm:schur}.
\item The descent pairs (\S\ref{ub-objects:sec:grids}): the normalized descent is a multiple of the
  transposed lift, and the parameters of the $j$-th pair correspond to $2\lambda(\tau)$.
  \lean{New.DescentPair.normalizedDescent_eq_smul_star_lift_swap}, \lean{New.Data.fitsParams_descentPair}.
  The text gives the statements only.
\item The description of the nodes of the grids $Y$ and $\tilde Y$ in
  Lemma~\ref{ub-objects:lem:grids}, their parameters and the rule for the signs.
  \lean{New.Data.gridSigns_eq_mul_sign}, \leanfile{New/UpperBound/Objects/Parameters}. The text
  proves the gap and the block rule, not this description.
\end{itemize}

\subsection*{Section~\ref{sec:der}: Derivatives of theta functions and the primes they count}
\begin{itemize}[leftmargin=2em,itemsep=1pt]
\item Step (3) of the sketch of proof of Theorem~\ref{der:thm:sharp}: the passage from
  $\mathrm{val}_6\le-\sum_\tau j_\tau$ to $S_\ell+2\sum_\tau v(\Pi_\tau)$ (choice of the $j_\tau$, Legendre's
  formula). \lean{New.Data.padicValNat_factorial_sub_min_le},
  \lean{New.Data.gammaPrime_mul_degree_sub_le_localPrimeSumTotal_add}. The text gives the result only.
\item In the same step: pairs of $b$-places sparse along the powers of $-\ell$ form a coupling family,
  and for $\ell>k$ the rest of the degree comes at valuation $\le0$.
  \lean{New.Data.isCouplingFamily_of_isNearPair}, \lean{New.Data.leastValLE_six_zero_of_inert}. The text
  gives the statements only.
\item Step (4) of the same sketch: the twelve copies are the jet series of a vector with integral
  values, and Fourier inversion passes to the indicator datum $\mathcal F_{\fm,\mathbf y_0}$.
  \lean{New.Data.exists_integral_isThetaFunction_twelveTaylor_eq},
  \lean{New.Data.exists_norm_le_norm_indicatorDatum}. The text gives a one-sentence summary.
\item The construction of the tight frames, of length $b'^2$ with $\ell\nmid b'$
  (Proposition~\ref{der:prop:product}, step (1)) and of length $M$ with $c_i\ge6fb$ (sketch of proof of
  Theorem~\ref{der:thm:sharp}). \lean{New.Data.exists_isTightFrame},
  \leanfile{New/Derivatives/Vacuum/ShortVectors}. The text uses their existence only.
\end{itemize}

\subsection*{Section~\ref{sec:ub-seesaw}: The seesaw identity}
\begin{itemize}[leftmargin=2em,itemsep=1pt]
\item The absorption of the unipotent parts of the two Bessel groups into the zero term (proof of
  Lemma~\ref{ub-seesaw:lem:double}).
  \lean{New.Seesaw.DoubleLiftData.exists_integral_constTermAtStringMap_mul_eq}. The text gives the
  mechanism as the documentation of the code describes it.
\item The finite computations in the proof of Lemma~\ref{ub-seesaw:lem:seesawDatum}: the evaluation of
  translates of the vacuum, the zero term as a sum over a neighbour lattice above $N$, and the
  denominator $\mathcal B^4$ of each $\mathrm{ev}_i$. \leanfile{New/UpperBound/Seesaw/AdmissibleEvaluation},
  \leanfile{New/UpperBound/Seesaw/NeighbourLattice}. The text gives the statements only.
\item The top-stage formula $\int_{[G(Y_m)]}F\bar v=C_l\langle v,\operatorname{res}_{Y_m}\psi\rangle$, used in
  the proof of Proposition~\ref{ub-seesaw:prop:chain}.
  \lean{New.Data.exists_integral_eq_mul_inner_restrictToFormY}. The text calls it a direct computation
  with the highest vector and the branching map.
\item In the proof of Lemma~\ref{ub-seesaw:lem:cut}: up to $R_{\mathrm f}(Y_j)$ the points
  satisfying~\eqref{ub-seesaw:eq:bounds} lie in finitely many cosets $cJ^\circ$, $c$ supported on
  $Q\cup\bar Q$. \lean{New.Grading.exists_finite_representatives}. The text gives the statement only.
\end{itemize}

\subsection*{Section~\ref{sec:ub-divis}: Divisibility and size of the composite lift}
\begin{itemize}[leftmargin=2em,itemsep=1pt]
\item The error terms in the proof of Lemma~\ref{ub-divis:lem:only}: $x^{60}$ per pair $(\tau,e)$, and
  the surplus $a/x_1$ for $\tau\in S$, $e=1$. \lean{New.Data.posPart_sub_weight_le_posPart_longRow},
  \lean{New.Data.localPrimeSum_add_le_mul_tsum_partnerTerm_add}. The text gives the comparison term by
  term and states the errors.
\item In Lemma~\ref{ub-divis:lem:orbits}: $m$ divides the order of $-\ell$, the residue fields of the
  primes above $\ell$, and a full mark exists. \lean{New.Data.IsExceptional.m_dvd_orderOf},
  \lean{New.Data.exists_isFullMark}. The text sketches the first in one sentence and states the others.
\item The identification of $\mathcal Y_{\mathfrak l}$ with an edge Hecke operator (proof of
  Proposition~\ref{ub-divis:prop:col}): edges, letter matrices, integral source forms as families of
  vectors. \leanfile{New/UpperBound/Divisibility/HeckeEdges}. The text gives a one-sentence summary.
\item The rewriting of the glued double lifts at $\Sigma_N$ as $c$ times a chain of direct lifts of
  $\mathrm{Up}^{m/2}\varphi$, with $\abs{\iota c}\le N^{mfb}$ (proof of Proposition~\ref{ub-divis:prop:N}).
  \lean{New.Data.doubleLiftChain_gluedDoubleLiftDatum_eq_smul_chain}, \lean{New.Data.norm_gluedFactor_le}.
  The text gives the statements only.
\item The colength in the same proof: the forms integral for $L_0$ have colength dividing $N^l$ in
  $\mathcal A^\vee$. \lean{New.Data.colengthDvd_integralFormsAt_diagPointN}. The text gives the
  statement only, not the comparison of the lattices class by class.
\item The count of at most $e^{x^{30}}$ classes at levels containing the principal level
  $\pi_M\pi_N$, used in the proof of Theorem~\ref{ub-divis:thm:size}.
  \lean{New.Data.card_stdClasses_le_exp}, \lean{New.Space.card_stdClasses_le}. The text gives the
  statement only.
\end{itemize}

\subsection*{Section~\ref{sec:ub-eldiv}: Elementary divisors of a Hecke operator}
\begin{itemize}[leftmargin=2em,itemsep=1pt]
\item The comparison of the weight of a monomial in one copy with the defect of its row, in the idea
  of proof of Theorem~\ref{ub-eldiv:thm:average}. \lean{New.FrobeniusCycle.abs_defect_sub_copyWeight_le}.
  The text attributes it to telescoping filtration levels; the choice of the weights, the cofactors
  at odd letters and the number $\nu/2$ of copies are not explained.
\item The weights of the auxiliary variable $\varsigma$ and the bookkeeping of the filtration in the
  proof of Theorem~\ref{ub-eldiv:thm:root}.
  \lean{New.FrobeniusCycle.RootAt.HasBoundedIncrements.aeval_mem_and_sub_mem_iteratedGaussFiltration}.
  The text writes ``filtering with a suitable weight''.
\item In the idea of the proof of Proposition~\ref{ub-eldiv:prop:count}: the integral shift into the
  cone, the passage from lattice points to volume as $j\to\infty$, and the comparison of flags and
  slab. \lean{New.FrobeniusCycle.exists_mem_defectCone_sum_le},
  \lean{New.FrobeniusCycle.exists_finrank_defectConditions_le_volume}. The text gives the statements
  only.
\item The volume of the cone in a slab, the part $\prod_{l<k}(q-1)/(q^{l+1}-1)$ (same proof).
  \lean{New.prod_mul_volume_real_inter_degreeSlab}. The text names the mechanism and the literature;
  the evaluation of the signed sum is not reproduced.
\item The construction of the points $z(\mathfrak y)$ and of the unipotent $g_{\mathfrak y}$ in the
  proof of Proposition~\ref{ub-eldiv:prop:deep}. \lean{New.FrobeniusCycle.entryCoeff},
  \lean{New.FrobeniusCycle.mem_lowPowerForms}. The text gives a summary.
\item That the $\varphi(g)^{\otimes\delta}$, $\delta+k\le q$, span the Schur algebra of a Levi subgroup
  over a local ring (proof of Proposition~\ref{ub-eldiv:prop:schur}).
  \lean{New.Average.eq_top_of_schurMap_eq_leviMonomialValues}. The text gives the statement, with the
  literature for the classical case.
\end{itemize}

\subsection*{Section~\ref{sec:sp-moment}: Spanning: the second moment, multiplicities and Bessel functions}
\begin{itemize}[leftmargin=2em,itemsep=1pt]
\item The bound for the pairing of $W^\flat$ with its translate over the open cell (paragraph
  \emph{The open cell}, \S\ref{sp-moment:ss:bessel}), used in Section~\ref{sec:sp-remain}.
  \lean{New.BesselSystem.enorm_integral_pairKernel_cut_le},
  \leanfile{New/Spanning/Bessel/OpenCellBound}. The text gives the statement and names its three
  ingredients.
\item Schur's bound for the Bessel system and its derivatives, in the proof of
  Theorem~\ref{sp-moment:thm:system}. \lean{New.IsBesselFunction.hasSchurBound_of_convolve_eq}. The
  text gives the statement and the mechanism in one sentence.
\item The proof of Proposition~\ref{sp-moment:prop:cut}: the covering argument for the decay
  integral and the majorant of the normalised Bessel average.
  \lean{New.DiscreteSeries.lintegral_decay_le_besselBallVolume},
  \lean{New.exists_isDecayMajorant_besselAverage}. The text gives a sketch.
\item The periods of discrete series (\S\ref{sp-moment:ss:appendix}): the proofs of
  Theorems~\ref{sp-moment:thm:laplace}, \ref{sp-moment:thm:stationary} and
  \ref{sp-moment:thm:restriction}. \lean{New.DiscreteSeries.IsLaplaceDatum.laplace_method},
  \lean{New.DiscreteSeries.IsSteep.exists_isStationary}. The text gives the statements and, for the
  first two, an outline as the documentation of the code describes it.
\end{itemize}

\subsection*{Section~\ref{sec:sp-unfold}: Unfolding the Poincar\'e series: the second moments}
\begin{itemize}[leftmargin=2em,itemsep=1pt]
\item The convergence of the amplitude sums and of the Poincar\'e series (after
  Definition~\ref{sp-unfold:def:amp}; Theorem~\ref{sp-unfold:thm:series}(a)), and the integral
  formula for $\mathrm{co}\varphi(c)\cdot W(x)$ in the proof of (b).
  \lean{New.Filling.tsum_amplitudeMajorant_lt_top},
  \lean{New.Filling.Fits.exists_dotProduct_besselFn_eq_integral}. The text gives one sentence for the
  first and the statement of the second.
\item The formula for the covolume $\nu$ and the identity $\int_{\mathsf
  U_\infty}\beta_c(uy)\,du=\nu$ (Definition~\ref{sp-unfold:def:amp}).
  \lean{New.Filling.lintegral_classCutoff_eq_unipotentCovolume}. The text describes $\nu$ only as an
  explicit product of powers and states the identity.
\item The verification of the standing hypotheses for the data $D$ (\S\ref{sp-unfold:ss:data}), among them the bound for
  the Siegel constant that \S\ref{sp-main:ss:proof} uses, and the size of $\mathcal
  D_{\mathrm{main}}/\mathcal D_{\mathrm{id}}$. \lean{New.Data.scaledSiegelConst_le},
  \lean{New.Data.abs_log_mainSecondMoment_div_identitySecondMoment_sub_le}. The text gives the
  statements and a two-sentence summary.
\item Proposition~\ref{sp-unfold:prop:cells}: in (a) the norm identity for $\Phi_v(\mathsf
  g(\gamma))$ and the criterion for $\mathsf g(\gamma)=0$, in (c) the decomposition $A=A_{\le
  j}\oplus(A\cap\gamma A)$. \lean{New.Filling.map_termDisc_eq_norm},
  \lean{New.Filling.termDisc_eq_zero_iff},
  \lean{New.Filling.CharsAgree.hasCellRank_of_add_finrank_eq}. The text gives the statements and
  names the mechanisms.
\item The values $\operatorname{sr}(u)=\tau\alpha$, $\operatorname{sr}(u')=\tau\alpha'$ of the
  Eichler transformations in the proof of Lemma~\ref{sp-unfold:lem:climb}.
  \lean{New.exists_mem_besselU_eq_eichlerMap}. The text gives the statement only.
\item The blocks of the Bessel subgroup (Proposition~\ref{sp-unfold:prop:sort}(b)) and the
  representation $r_\psi$ (Lemma~\ref{sp-unfold:lem:rpsi}(b), (c)).
  \lean{New.Filling.amplitudeSum_besselS_one_eq_ite},
  \lean{New.Filling.HasAmplitudeBlocks.adjoint_comp_rep_mul_left}. The text gives a sketch.
\item The bound $\sum_s\abs{\psi^\natural(s)}^2\le2$ in the proof of Theorem~\ref{sp-unfold:thm:good},
  through a local mass integral of Section~\ref{sec:sp-frames}.
  \lean{New.Filling.sum_norm_reducedCharMean_sq_le}. The text gives the statement and names the
  source of the bound.
\end{itemize}

\subsection*{Section~\ref{sec:sp-main}: The main term of the second moment}
\begin{itemize}[leftmargin=2em,itemsep=1pt]
\item In the proof of Theorem~\ref{sp-main:thm:sieve}: the count of frequencies behind the sum of the
  decay over nonzero frequencies, and the bound for $(\sum\Lambda_i)^2$.
  \lean{New.sum_height_inv_cube_le}, \lean{New.Filling.sq_sum_globalWeight_le}. The text gives the
  statements and one sentence for the second.
\item That $E_{\mathcal P}$ is a finite place in the sense of Section~\ref{sec:sp-frames} with a
  large lattice datum of modulus $M$ (\S\ref{sp-main:ss:local}). \lean{New.Filling.exists_isFillingData},
  \lean{New.Filling.exists_integralPlace}. The text gives the statement only.
\item The constants: $c_{\mathrm n}$ and the identity for $\rho_\alpha\mathsf h_\varrho(\alpha,1)$, the
  bound for $\operatorname{mr}(Z)$ in step (iv) of the proof of Proposition~\ref{sp-main:prop:siegel},
  the value of $\prod_iP_i$ and the final comparison with $1+b^{-4}$ (\S\ref{sp-main:ss:proof}).
  \lean{New.Filling.IsFillingData.exists_scaledSiegelConst_eq_and_massRatio_le},
  \lean{New.Filling.sieve_constant_le}. The text gives the statements only.
\end{itemize}

\subsection*{Section~\ref{sec:sp-frames}: Two frames at one place: the local estimates}
\begin{itemize}[leftmargin=2em,itemsep=1pt]
\item Proposition~\ref{sp-frames:prop:gains}(i): solving the equation of $Q_j$ in two blocks of
  orthonormal coordinates, and the Fourier kernel between two quadrics.
  \lean{New.Floors.Gauge.fourierSqNormLe_floorQuadricMeasure}. The text gives one sentence.
\item Proposition~\ref{sp-frames:prop:gains}(iii), (iv): the bound $\int\delta_a^{\mathsf
  s}\le\Lambda^{10^4n}$ on each floor, and the floor-by-floor induction with its constants.
  \lean{New.Floors.Gauge.floorMassLe_depth_pow},
  \lean{New.Floors.Gauge.lintegral_weightedSquareIntegral_le}. The text gives the mechanism as the
  documentation of the code describes it, without the constants.
\item The shell sum (paragraph \emph{Shells}, \S\ref{sp-frames:ss:gains}): the covering of a shell by
  position balls, and why bounded depths bound $(y-1)/M$ (comparison of gauges at two points).
  \lean{New.Floors.Gauge.hasMatrixSize_sub_one_of_stableProducts},
  \lean{New.Floors.Size.GaugesDifferLe.of_succ}. The text gives the statements and an outline.
\item The comparison $c_{\mathrm H}\vol\U(k)\le\vol\U(r)$ at $\C$ (\S\ref{sp-frames:ss:real}).
  \lean{New.Floors.RealPlaceSizes.haarRatio_ne_zero_and_groupMeasure_eq_smul}. The text gives the
  statement only.
\item The proof of Theorem~\ref{sp-frames:thm:findecay}: the gauge adapted to two heads, the common
  divisor $\eta$, and the exponent arithmetic. \leanfile{New/Spanning/Frames/FiniteDepth},
  \leanfile{New/Spanning/Frames/FiniteShells}. The text gives a sketch.
\item The top floor (Proposition~\ref{sp-frames:prop:top}; step (iii) of the proof of
  Proposition~\ref{sp-main:prop:siegel}): local constancy and compact support of the scaled orbital
  integrals, so that $\Phi$ is a step function. \leanfile{New/Spanning/Frames/TopCorrelation}. The
  text gives the statement only.
\item The passage to several places in the proof of Theorem~\ref{sp-frames:thm:fixed}, as products
  over the fibres $\{w,\bar w\}$. \lean{New.Filling.IsFillingData.massIntegral_pi_mul_le}. The text
  gives one sentence.
\end{itemize}

\subsection*{Section~\ref{sec:sp-remain}: The remaining terms of the second moment}
\begin{itemize}[leftmargin=2em,itemsep=1pt]
\item Schur's test in the proof of Proposition~\ref{sp-remain:prop:size}: the majorant of the weight
  $\prod_v\nu_v(y)^{-(k+i)}$ behind the bound for $\int\varphi$, and the packing of the isotropic
  parts at a cost $\tau$. \lean{New.RealPlace.exists_majorant_volumeRatio_inv_pow},
  \lean{New.Filling.prod_kernelBound_le}. The text gives the statements only.
\item The proof of Proposition~\ref{sp-remain:prop:gain}: the discriminant of $L^+_\gamma$, and the
  lower bound $\operatorname{mn}(s)$ for the minor of the two $z$-forms in the cut.
  \lean{New.Filling.exists_hasLevelLifts_linkDiscSq_le},
  \lean{New.Filling.HasIntegralLifts.le_genIndex_det_zRowBlockDiagonal}. The text gives one sentence
  for each.
\item The bound $2^k$ for the zeta function of sublattices in the proof of
  Proposition~\ref{sp-remain:prop:lattices}. \lean{New.Filling.tsum_index_saturatedLattice_le}. The
  text gives the statement only.
\item The proof of Proposition~\ref{sp-remain:prop:frames}: the power of $b$ that each of the three
  rounds yields, and the exponent $2000n^6(l+1)f$.
  \leanfile{New/Spanning/RemainingTerms/FramesOverLattice}. The text gives the three rounds as the
  documentation of the code describes them.
\end{itemize}

\subsection*{Section~\ref{sec:det-chain}: The lower bound: Theorem C and the determinant of a chain of lifts}
\begin{itemize}[leftmargin=2em,itemsep=1pt]
\item The conversion of the finite square law of the two composites into the commutation
  relation~(i) of a triangle, by linear algebra over the residue ring (proof of
  Theorem~\ref{det-chain:thm:triangle}). \leanfile{New/Determinant/Chain/SquareLaw},
  \lean{New.Data.exists_triangleLift_triangleTransposedLift_comm}. The text gives a one-sentence
  summary.
\item The adjoint law in the squares: the second commutation relation, for the transposed lifts
  along the reversed pairs, with constants of modulus~$1$ (proof of
  Theorem~\ref{det-chain:thm:triangle}). \lean{New.Data.exists_triangleLift_triangleReversedLift_comm},
  \lean{New.Data.exists_triangleTransposedLift_adjoint_comm}. The text gives a one-sentence summary.
\item The index $[K_0:J_\vartheta(Y_m)]=C'(\Nm Q)^{u}\delta^mw_m/(w_0w_1^m)$, by a count of block
  triangular groups (proof of Theorem~\ref{det-chain:thm:norm}).
  \lean{New.Data.exists_stdLevelIndex_lastLevel_eq}. The text gives the statement only.
\item The spanning of the principal forms on $Y_0$ by forms with short lifts (proof of
  Theorem~\ref{det-chain:thm:aux}). \lean{New.Data.exists_span_short_lifts_eq_top}. The text gives
  the route (first path, $m^2$ turns, last path, delta forms), not the estimates by which deep-level
  forms on $\tilde Y^0_0$ detect the principal forms on $Y_0$
  (\lean{New.Data.exists_eq_zero_of_orthogonal_gridLevelForms}), nor the diagonal turns
  (\lean{New.Data.exists_transfersPeriods_first_last}).
\item The integrality of $n_\Xi\Xi$ away from $\Nm Q$, with $\log n_\Xi\le x^{99}b^2$, and the bound
  for the conjugates of the period at $\vartheta=1$, on the last path (proof of
  Theorem~\ref{det-chain:thm:aux}). \lean{New.Data.exists_mul_lastPathPeriod_mem_embValRing},
  \lean{New.Data.norm_lastPathPeriod_le}. The text gives the statements only.
\end{itemize}

\subsection*{Section~\ref{sec:det-exchange}: The exchange identity}
\begin{itemize}[leftmargin=2em,itemsep=1pt]
\item The reduction in step (ii) of the proof of Proposition~\ref{det-exchange:prop:top}, from the
  coefficients at $\Xi$ unimodular at $S$ to the points where $L_0^*\Xi L_0$ is unimodular at all
  finite places, through the partial means along the extra plane of $\widetilde W$ and
  Fact~\ref{det-exchange:fact:lib}(b).
  \lean{New.Exchange.IsEquivariantLift.exists_linear_partialMean},
  \lean{New.Exchange.IsEquivariantLift.shearInvariant_stdColDatum}. The text states the reduction
  and names its inputs; it does not show how the partial means meet the hypotheses of the criterion.
\item The constant $C_\nu$ of step (iii) of the proof of Proposition~\ref{det-exchange:prop:top},
  with $\abs{C_\nu}=C_2\abs{I_0(\nu)}$, and the period identity for the vectors of step (iv).
  \lean{New.Exchange.exists_integral_thetaTerm_eq_of_periods},
  \lean{New.Exchange.IsWittBump.isPeriodIdentity}. The text gives the shape of the unfolding and
  the definition of the vectors; the constants of the unfolding are not given.
\item The two analytic inputs of the non-vanishing in the proof of
  Proposition~\ref{det-exchange:prop:equalrank}: the vanishing of the large side in case (b) for
  $V'=J^-$, and the vector $v'$ whose cone integral has prescribed theta integrals up to a positive
  analytic factor $B(y)$. \lean{New.Exchange.centreCoeff_eq_zero_of_not_representedBy},
  \lean{New.Exchange.exists_thetaLarge_cone_eq_mul}. The text gives the mechanism of each in one
  sentence.
\item The archimedean step of the proof of Theorem~\ref{det-exchange:thm:explicit} (circle mean,
  average over $\U(I)(\R)$ and polarisation). \lean{New.Rows.exists_mul_apply_eq_apply_matchedPoly}.
  The text gives a one-sentence summary.
\end{itemize}

\subsection*{Section~\ref{sec:ol-values}: One lift: its determinant, false lifts and Hecke L-values}
\begin{itemize}[leftmargin=2em,itemsep=1pt]
\item At the bottom of the tower in the proof of Theorem~\ref{ol-values:thm:periods}: the
  approximation of general data by indicators of small cosets, and the form of the mean of their
  translates under the level group. \lean{New.IsSchwartzBruhat.mem_of_eventually_coset},
  \lean{New.DescentPair.eq_smul_translatedBlockFn}. The text gives the mechanism as the
  documentation of the code describes it.
\item In the step of the tower (proof of Theorem~\ref{ol-values:thm:periods}): the boundedness of
  the theta kernel of $(W',I_n)$ minus its zero term after a rational translation, a reduction-theory
  argument, and the inequalities on the size of the Hecke root.
  \lean{New.SiegelWeil.exists_norm_theta_sub_zeroTerm_le},
  \lean{New.SiegelWeil.lift_eq_zero_of_heckePolyOp_eq_zero}. The text gives the statements only.
\item The seesaw at one point, which factors the pairing of the test form with the theta kernel of
  $p$ along the blocks of sizes $12,4,4,2z-20$ (proof of Proposition~\ref{ol-values:prop:test}).
  \lean{New.Data.inner_pointForm_kernelForm}. The text gives the statement only.
\item The integrality of the integral algebraic part: $n_L$ clears the denominators of the kernel
  values and of the stabilisers (proof of Proposition~\ref{ol-values:prop:untwisted}).
  \lean{New.Data.CriticalPair.norm_integralAlgebraicPart_one_le_one},
  \lean{New.Data.CriticalPair.weylGauge_denomBound_smul_pointKernel_le_one}. The text gives the
  statement only.
\item Results (2) and (3) at the end of the section: the bound for the product of the projections
  of two test forms to the line of $\varphi^\circ$, through the central action of $\U(1)$, and the
  factor gained by the period of the determinant theta kernel when a prime is added to the modulus.
  \lean{New.Selmer.DenominatorLe.norm_inner_starProjection_le}, \lean{New.Rows.detPeriod_insert}.
  The text gives the statements only.
\end{itemize}

\subsection*{Section~\ref{sec:ol-compare}: One lift: the comparison operator}
\begin{itemize}[leftmargin=2em,itemsep=1pt]
\item The construction of the glued lattice $\gamma_\Lambda\cO^n$ and the proof that it is glued
  (\S\ref{ol-compare:sec:gluing}). \lean{New.Rows.DividedSplitting.exists_isGlued},
  \lean{New.Rows.DividedSplitting.glueMatrix}. The text gives the statement only.
\item The size estimates in Propositions~\ref{ol-compare:prop:first}(b)
  and~\ref{ol-compare:prop:glue}(b). \lean{New.Data.firstFactorSize_sq_le_exp},
  \lean{New.Rows.M_pow_denominatorExponent_sq_mul_le_exp}. The text gives the ingredients, not the
  arithmetic.
\item The holomorphy of the $\mathsf F_i$, and the dominance of the weights of $\bar y$,
  $\theta^\circ$ and $\mathsf F_i$ with gaps $300n_*$ (proof of Theorem~\ref{ol-compare:thm:span}).
  \lean{New.Data.isAdelicHol_choiceFn}, \lean{New.Data.conjHolomorphicWeight_eq_add_kernelWeight}.
  The text gives the statement only.
\item The analytic estimates of steps 3 and 4 of the proof of Theorem~\ref{ol-compare:thm:span}:
  the approximation of $c_P$ by a kernel with a finite datum, and the isolation estimate by a mean
  over a circle of peaks. \lean{New.Data.exists_datum_norm_modKernel_sub_polyMatrixCoeff_le},
  \lean{New.Data.norm_apply_one_le_of_forall_isPeak}. The text gives the mechanism of each in one
  sentence.
\end{itemize}

\subsection*{Section~\ref{sec:cg-selmer}: Congruences: regular families of twists, Galois representations and Selmer groups}
\begin{itemize}[leftmargin=2em,itemsep=1pt]
\item The Burnside-type step in the proof of Lemma~\ref{cg-selmer:lem:typerow}: from an eigenvalue
  of $R(g)$ with $r'$ Frobenius conjugates distinct modulo $\ell$, every lattice spans
  $\mathrm M_{r'}(\cO)$ and $\phi^{r'}$ acts as a scalar.
  \lean{New.Selmer.IsLatticeIn.span_range_eq_top_of_map_pow_eq_map_conj}. The text gives the
  statement only.
\item The integrality of the measures $F$, and the identity between $\abs{\Pi_\vartheta}^2$ and
  $\Nm Q_\vartheta\abs{L(1,(\nu\vartheta)^{-1})}^2$, in the idea of proof of
  Proposition~\ref{cg-selmer:prop:stickel}.
  \lean{New.Data.CriticalPair.isIntegralMeasure_integralAlgebraicPart},
  \lean{New.Data.CriticalPair.exists_isIntegralMeasure_norm_sq_detPeriod},
  \lean{New.Data.CriticalPair.exists_mul_norm_sq_detPeriod_eq}. The text gives the statements only.
\item Sah's lemma ($z(H)$ spans $\Lambda/I$) and the tilting lemma (at most $2n$ residue classes of
  slopes are excluded at each level), in the idea of proof of Theorem~\ref{cg-selmer:thm:descent}.
  \lean{New.Selmer.SelmerDatum.span_image_sahGroup_eq_top}, \lean{New.Selmer.tilting}. The text
  gives the statements only.
\item The sharp polynomial in \S\ref{cg-selmer:sec:proof}: the Selmer modules modulo
  $\ell^{\mathsf M}$ of all $\Lambda\psi_{\mathsf x}$ are the kernels of one matrix over $\cO[X]$,
  and some minor is sharp at a given point. \lean{New.Selmer.exists_sharp_polynomial}. The text
  gives the statement in one sentence.
\item The base data for patching of a lift setting (end of \S\ref{cg-selmer:sec:galois}): the
  construction of the forms and the bookkeeping of product labels and parahoric types.
  \lean{New.Selmer.exists_baseLevel}, in
  \leanfile{New/Congruences/GaloisRepresentations/BaseLevelOfLabel}. The text states what the
  result gives, not how the data are built.
\end{itemize}

\subsection*{Section~\ref{sec:cg-first}: Congruences to first order: Ribet's method and the bound for denominators}
\begin{itemize}[leftmargin=2em,itemsep=1pt]
\item The relation $\mathrm{red}(c\,\partial_j(g-1))=\mathrm{red}(\mathrm{Nm}_jZ_g)$ with
  $\delta'\mathrm{Nm}\,Z_g=0$, the source of the last two properties of
  Theorem~\ref{cg-first:thm:tie} (end of its idea of proof).
  \lean{New.Patching.exists_map_nsmul_mul_sub_one_eq_sum_mul},
  \lean{New.Patching.BaseLevel.heckeAction_mul_norm_mul_eq_zero}. The text names the two inputs,
  Fact~\ref{cg-first:fact:chain}(ii) and Fact~\ref{cg-first:fact:tw} for $\Pi$, but does not carry
  out the computation.
\item The congruence used in the proof of Proposition~\ref{cg-first:prop:one}.
  \lean{New.Selmer.exists_prod_sub_det_one_sub_eq_zero}. The text gives the statement only.
\end{itemize}

\subsection*{Section~\ref{sec:cg-lifting}: The local loss above \texorpdfstring{$\ell$}{l}: crystalline classes in large weight}
\begin{itemize}[leftmargin=2em,itemsep=1pt]
\item The stop at a long gap, in the proof of Theorem~\ref{cg-lifting:thm:cuts}: how the point of
  Lemma~\ref{cg-lifting:lem:costs}(c) gives a stop of depth $2b$ at the pivot that obeys the rule
  of Lemma~\ref{cg-lifting:lem:pairing}(b).
  \lean{New.Lift.KummerChart.PairingSide.pointsOver_longGapDepth_nonempty},
  \lean{New.Lift.KummerChart.PairingSide.ruleBound_le_inf_add_longGapDepth}. The text gives the
  statement only.
\item In the proof of Theorem~\ref{cg-lifting:thm:generic}, why the depth must be raised at
  level~$1$ at the letters with a negative far window.
  \lean{New.Lift.KummerChart.PairingSide.negativeWindowDepth},
  \lean{New.Lift.KummerChart.PairingSide.isSafeMoveFst_of_negativeWindowDepth}. The text gives what
  this step costs, not why it is needed.
\item In the same proof, the step by which, among $f_\Gamma\cdot\mathrm{mult}+1$ points, some
  primitive combination gains one step of order.
  \lean{New.Lift.KummerChart.JumpLattice.exists_sum_smul_mem_locKisinPart_sup}. The text gives the
  statement only, as the documentation of the code states it.
\item Lemma~\ref{cg-lifting:lem:local}, on the local lattices of a conjugate self-dual
  representation. \lean{New.Selmer.IsConjSelfDualOfLevel.hasJumps_localLattice},
  \lean{New.Selmer.IsConjSelfDualOfLevel.inClass_localLattice},
  \lean{New.Selmer.isTameType_of_roots}. The text gives the statement only.
\end{itemize}

\endgroup
